% =============================================================================
% PROPOSED CHANGES to _ICML_2026__Empirical_GPs
%
% Standalone, compilable preview of every edit we propose. Nothing here modifies
% main.tex -- read the PDF, then move the marked blocks across.
%
% Build:  cd paper && pdflatex paper_changes.tex && pdflatex paper_changes.tex
%
% Notation matches main.tex exactly: K training functions, C_i posterior
% covariances, \tilde m_i mean adjustments, W = K_{*X} K_{XX}^{-1},
% delta_Sigma = Sigma - K_{XX}.
% =============================================================================
\documentclass[10pt]{article}
% NOTE: compiled with plain `article` because icml2026.sty needs forloop.sty,
% which is not installed here. Text and math are unaffected; only the ICML
% chrome differs. When pasting into main.tex nothing below changes.
\usepackage[margin=1in]{geometry}

\usepackage{amsmath,amssymb,amsthm,bm}
\usepackage{booktabs}
\usepackage{xcolor}
\usepackage[colorlinks=true,allcolors=blue]{hyperref}
% Preview-only shims: cleveref/mdframed are not installed here. The paper HAS
% cleveref, so \Cref is kept verbatim in the pasteable text below.
\providecommand{\Cref}[1]{Appendix~\ref{#1}}

% --- macros borrowed from main.tex so the blocks drop in verbatim ------------
\newcommand{\vx}{\mathbf{x}}
\newcommand{\bX}{\mathbf{X}}
\newcommand{\bZ}{\mathbf{Z}}
\newcommand{\bSigma}{\boldsymbol{\Sigma}}
\newcommand{\domain}{\mathcal{X}}

% --- presentational only: NOT to be copied into the paper -------------------
\definecolor{addgreen}{RGB}{0,110,0}
\newenvironment{proposed}[1]
  {\par\medskip\noindent{\color{addgreen}\rule{\linewidth}{1pt}}\par
   \nopagebreak{\color{addgreen}\bfseries #1}\par\nopagebreak\smallskip
   \begin{quote}\itshape}
  {\end{quote}{\color{addgreen}\rule{\linewidth}{0.4pt}}\par\medskip}
\newcommand{\editnote}[1]{\par\smallskip{\color{addgreen}\small\textbf{Where:} #1}\par\smallskip}

\title{Proposed Changes to \emph{Empirical Gaussian Processes}\\
\large A standalone preview --- not part of the paper}
\author{Draft for review}
\date{\today}
\begin{document}
\maketitle

\section*{Summary of proposed changes}

Four technical components exist in the implementation but are absent from the
manuscript. All four are drafted below as \textbf{appendix} material, so the main
text grows by only three sentences.

\begin{enumerate}\itemsep2pt
\item \textbf{Covariance shrinkage in the M-step} (\Cref{apd:shrinkage}). The
  manuscript's Conclusion currently lists this as an open limitation
  (``\emph{possibly requiring shrinkage regularization such as an
  Inverse-Wishart prior}''), and the derivation is commented out around
  \texttt{main.tex:647}. We show the trace-matched blend is \emph{exactly} MAP-EM
  under that Inverse-Wishart prior, which turns the limitation into a result.
\item \textbf{Base-augmented conditioning} (\Cref{apd:base_augmented}).
  \textbf{This one is a reproducibility fix, not an enhancement:} the 7D numbers
  in \texttt{tab:hd\_nll} are produced \emph{with} a fitted additive base kernel
  that the manuscript never mentions. A reader instantiating the plain model will
  not reproduce that table.
\item \textbf{Multi-output Empirical GP} (\Cref{apd:multioutput}) --- dense,
  non-separable cross-output covariance.
\item \textbf{Multi-task Empirical GP} (\Cref{apd:multitask}) --- long-format,
  partial observation, heterogeneous per-task domains.
\end{enumerate}

Section~\ref{sec:mainedits} lists the three main-text sentences; Appendix~\ref{apd:variants}
is the appendix section to paste in wholesale.

\section{Main-text edits (three sentences plus one rewrite)}
\label{sec:mainedits}

\begin{proposed}{Edit 1 --- add after the EM derivation in Section 3}
When $K \le M$ the maximum-likelihood covariance is rank-deficient; we regularize
the M-step by shrinking toward a structured target, which is exactly MAP-EM under
an Inverse-Wishart prior (\Cref{apd:shrinkage}).
\editnote{Section 3, immediately after the M-step update.}
\end{proposed}

\begin{proposed}{Edit 2 --- add where conditioning is introduced}
At conditioning time the empirical covariance may be augmented with a fitted
additive base kernel, which guarantees the model is never asymptotically worse
than that kernel (\Cref{apd:base_augmented}).
\editnote{Section 3, or Section 4.5 with the continuous-domain variant.}
\end{proposed}

\begin{proposed}{Edit 3 --- REQUIRED: name the variant behind Table~2}
\ldots we evaluate our continuous-domain EM-EGP variant \textbf{with a fitted
additive base kernel} (\Cref{apd:base_augmented}) on a 7-dimensional
hyperparameter optimization task from LCBench.
\editnote{Section 5.5, in the sentence introducing \texttt{tab:hd\_nll}. Without
this the table is not reproducible.}
\end{proposed}

\begin{proposed}{Edit 4 --- rewrite one Conclusion sentence}
\textbf{Currently:} ``With few historical realizations, the empirical covariance
is high-variance and can be singular, \emph{possibly requiring shrinkage
regularization such as an Inverse-Wishart prior}.''

\smallskip
\textbf{Proposed:} ``With few historical realizations, the empirical covariance is
high-variance and can be singular; we address this with a trace-matched shrinkage
M-step that is exactly MAP-EM under an Inverse-Wishart prior
(\Cref{apd:shrinkage}).''
\editnote{Conclusion, limitations paragraph. This stops us listing as future work
something the paper now does.}
\end{proposed}

\begin{proposed}{Edit 5 --- OPTIONAL, one sentence for the model family}
The same estimator extends to vector-valued and partially-observed multi-task
settings (\Cref{apd:multioutput} and \Cref{apd:multitask}).
\editnote{Section 3, or Related Work.}
\end{proposed}

\appendix

\section{Model Variants and Regularization}
\label{apd:variants}

\emph{(Paste this whole section into the appendix, after ``Additional Empirical
Evaluations''.)}

\subsection{Regularized M-step via Covariance Shrinkage}
\label{apd:shrinkage}

The M-step update
\begin{equation}
\label{eq:sigma_ml}
    \bSigma_{\text{ML}}
    = \frac{1}{K-1}\sum_{i=1}^{K}\Big[\mathbf{C}_i
      + \tilde{\mathbf{m}}_i \tilde{\mathbf{m}}_i^{\top}\Big]
\end{equation}
is rank-deficient whenever $K \le M$: the outer products contribute at most $K-1$
directions, and the posterior terms $\mathbf{C}_i$ vanish for completely observed
functions. The estimator is therefore singular precisely in the regime the method
targets --- many reference locations $\bZ$, few historical realizations.

\paragraph{Trace-matched shrinkage.}
We regularize the M-step by blending toward a structured target $\mathbf{B}$ (in
practice the base-kernel gram $\mathbf{K}_{\bZ\bZ}$),
\begin{equation}
\label{eq:shrinkage}
    \bSigma_\alpha = (1-\alpha)\,\bSigma_{\text{ML}} + \alpha\, s\, \mathbf{B},
    \qquad
    s = \frac{\operatorname{tr}(\bSigma_{\text{ML}})}{\operatorname{tr}(\mathbf{B})},
    \qquad \alpha \in [0,1].
\end{equation}
Trace matching through $s$ makes the blend scale-free: only the \emph{correlation
structure} of $\mathbf{B}$ is imposed, never its magnitude, so $\alpha$ carries the
same meaning across datasets with different output scales.

\paragraph{Justification as MAP-EM.}
Equation~\eqref{eq:shrinkage} is not a heuristic. Placing an Inverse-Wishart prior
$\bSigma\sim\mathcal{IW}(\boldsymbol{\Psi},\nu)$ on the covariance, the M-step
maximizer of the penalized expected complete-data log-likelihood is
\begin{equation}
    \bSigma_{\text{MAP}}
    = \frac{\boldsymbol{\Psi} + K\,\bSigma_{\text{ML}}}{K + \nu + M + 1}.
\end{equation}
Choosing the prior scale $\boldsymbol{\Psi} = (\nu + M + 1)\,s\,\mathbf{B}$ recovers
\eqref{eq:shrinkage} exactly, with
\begin{equation}
\label{eq:alpha_iw}
    \alpha = \frac{\nu + M + 1}{K + \nu + M + 1}.
\end{equation}
So $\alpha$ is an interpretable prior strength, and \eqref{eq:alpha_iw} shows the
regularization decays as $O(1/K)$: the estimator is automatically less regularized
as historical data accumulates, and $\alpha\to 0$ recovers maximum likelihood.
Because $s$ is re-matched to the data each iteration, the procedure is
empirical-Bayes rather than fully Bayesian.

\paragraph{Why a convex blend rather than an additive one.}
The additive alternative $\bSigma_{\text{ML}} + \alpha\mathbf{B}$ is equally natural
but is \emph{not} trace-bounded: every M-step adds $\alpha\operatorname{tr}(\mathbf{B})$,
so the iterates can grow without bound and the EM map need not admit a fixed point.
The convex blend \eqref{eq:shrinkage} maps the compact convex set
$\{\bSigma\succeq 0 : \operatorname{tr}(\bSigma)\le T\}$ into itself, so a fixed
point exists by Brouwer's theorem. We therefore use the convex form during
\emph{estimation}, and reserve additive augmentation for \emph{conditioning}
(\Cref{apd:base_augmented}), where the target data --- not the EM recursion ---
controls the added magnitude.

\paragraph{Empirical behavior.}
Table~\ref{tab:shrinkage_spectrum} reports the effect on the LCBench corpus ($M=200$
reference points, $K=25$ historical tasks). The unregularized estimator is severely
rank-deficient: its effective rank --- the exponential of the spectral entropy ---
is $2.89$, with $99.98\%$ of the trace in the leading $22$ directions and a median
eigenvalue of $2.0\times10^{-4}$, below the conditioning noise floor. Shrinkage
restores conditioning and repairs the resulting overconfidence, raising $95\%$
predictive coverage from $0.65$ to $0.98$ and reducing predictive NLL by about
$8$ nats.

\begin{table}[ht]
\small\centering
\caption{Effect of M-step shrinkage on the empirical covariance spectrum and on
predictive calibration (LCBench, $M=200$, $K=25$). Coverage and NLL are measured at
acquired points for the frozen-prior variant.}
\label{tab:shrinkage_spectrum}
\begin{tabular}{lcc}
\toprule
 & $\alpha = 0$ & $\alpha = 0.3$ \\
\midrule
Effective rank (of $M=200$) & 2.89 & 5.05 \\
Trace in top 22 modes       & 0.9998 & 0.9624 \\
Median eigenvalue           & $2.0\times10^{-4}$ & $1.6\times10^{-2}$ \\
\midrule
$95\%$ coverage             & 0.65 & \textbf{0.98} \\
Predictive NLL              & 8.30 & \textbf{0.10} \\
\bottomrule
\end{tabular}
\end{table}

Better calibration does not, however, imply better decision-making. When the target
task is drawn from the same distribution as the historical corpus, shrinkage
\emph{degrades} downstream Bayesian-optimization regret even as it improves
calibration, because the empirical covariance is already well matched and the blend
only dilutes it. Under distribution shift the sign reverses and shrinkage helps.
Shrinkage and target adaptation are therefore substitutes rather than complements,
and $\alpha$ is best read as a transfer-distance dial: $\alpha = 0$
in-distribution, $\alpha > 0$ as the target moves away from the corpus.

\subsection{Base-Augmented Conditioning}
\label{apd:base_augmented}

At conditioning time we augment the empirical covariance with a parametric base
kernel,
\begin{equation}
\label{eq:base_augmented}
    k(\vx,\vx') = k_{\text{emp}}(\vx,\vx') + k_{\theta}(\vx,\vx'),
\end{equation}
fitting $\theta$ and the likelihood noise by maximizing the marginal likelihood of
the \emph{target} task, with the empirical mean and covariance held fixed. Unlike a
convex combination, each component keeps its own magnitude. This gives graceful
degradation: when the empirical prior is informative the fitted outputscale of
$k_\theta$ stays small and the model behaves like the Empirical GP; when it is not,
$k_\theta$ grows to explain the target data and the model reverts to a standard GP
with that base kernel. The augmented model is therefore never asymptotically worse
than its base kernel --- a guarantee the raw empirical prior cannot provide at
finite rank.

Which parameters are free is a modeling choice with a clear reading. Fitting all of
$\theta$ suits single-task surrogates; freezing the lengthscales of a pre-fit base
kernel and fitting only its outputscale suits transfer, where learned cross-task
structure should be reused rather than re-estimated from a handful of target
points. We find the plain additive form \eqref{eq:base_augmented} outperforms richer
parameterizations that additionally free the empirical component's magnitude, which
over-parameterize a small target design.

\paragraph{Off-grid augmentation.}
Care is required in the continuous-domain variant. Recall the interpolated
covariance
$\bSigma(\bX_*,\bX_*) = \mathbf{A} + \mathbf{W}\bSigma\mathbf{W}^{\top}$ with
$\mathbf{W} = \mathbf{K}_{*X}\mathbf{K}_{XX}^{-1}$ and Nystr\"om residual
$\mathbf{A} = \mathbf{K}_{**} - \mathbf{K}_{*X}\mathbf{K}_{XX}^{-1}\mathbf{K}_{X*}$.
The augmentation must add the \emph{full} $k_\theta(\bX_*,\bX_*)$ at the query
points, not the projection $\mathbf{W}k_\theta(\bX,\bX)\mathbf{W}^{\top}$ of the
base kernel through the reference set; the latter discards exactly the residual
variance $\mathbf{A}$ that makes the augmentation useful away from the reference
grid. On-grid ($\bX_*\subseteq\bZ$) the two coincide, which is why the distinction
appears only in the continuous-domain setting.

\subsection{Multi-Output Empirical Gaussian Processes}
\label{apd:multioutput}

The construction extends to vector-valued $f:\domain\to\mathbb{R}^{m}$ with no extra
modeling assumptions, because the empirical estimator is indifferent to how the
index set is structured. Given $K$ historical realizations observed on a shared grid
of $n$ points for all $m$ outputs, we vectorize each realization,
$\operatorname{vec}(\mathbf{Y}_i)\in\mathbb{R}^{nm}$, and form the empirical mean and
covariance in the vectorized space. The resulting
\begin{equation}
    \bSigma^{\text{MO}} \in \mathbb{R}^{nm\times nm},
    \qquad
    \big[\bSigma^{\text{MO}}\big]_{(j,a),(l,b)}
    = \widehat{\operatorname{Cov}}\!\left[f^{(a)}(x_j),\, f^{(b)}(x_l)\right]
\end{equation}
is a \emph{dense} cross-output covariance, capturing correlations between different
outputs at different inputs directly from data rather than imposing a separable
intrinsic-coregionalization structure
$\bSigma^{\text{MO}} = \mathbf{B}\otimes\mathbf{K}$. Separability is a strong
assumption --- it forces every output to share one input correlation pattern --- and
the empirical construction does not need it. Predictions are returned as a multitask
multivariate normal over the $nm$ latent values.

Since $\bSigma^{\text{MO}}$ has rank at most $K$, when $K < nm$ we factor the
vectorized realizations by SVD and propagate the low-rank factors, reducing
conditioning cost from $O((nm)^3)$ to $O(K^2 nm)$ without approximation.

\subsection{Multi-Task Empirical Gaussian Processes}
\label{apd:multitask}

Appendix~\ref{apd:multioutput} assumes every output is observed at every input. Many
transfer settings violate this: tasks are observed at different locations, over
different domains, and in different numbers. We therefore also provide a
\emph{long-format} variant in which an input is a pair $(x,t)$ with
$t\in\{1,\dots,T\}$ a task index, so each observation carries its own task label and
no completeness assumption is needed.

Historical data may be heterogeneous across tasks: task $t$ contributes curves on its
own grid $\bX^{(t)}\in\mathbb{R}^{n_t}$, with $n_t$ and the domain differing freely
between tasks. The only requirement is that the $K$ realizations are \emph{paired}
across tasks --- realization $i$ of task $t$ and realization $i$ of task $t'$ arise
from the same underlying draw --- since this pairing identifies the cross-task
covariance
\begin{equation}
    \widehat{\operatorname{Cov}}\!\left[f^{(t)}(x),\, f^{(t')}(x')\right]
    = \frac{1}{K}\sum_{i=1}^{K}
      \big(f^{(t)}_i(x) - \hat\mu^{(t)}(x)\big)
      \big(f^{(t')}_i(x') - \hat\mu^{(t')}(x')\big),
\end{equation}
evaluated by interpolating each task's empirical basis to the queried locations.
Within-task covariance is the one-dimensional empirical kernel of the main text; the
cross-task blocks are estimated the same way, so a single estimator supplies the
entire $\big(\sum_t n_t\big)$-dimensional joint covariance. This is the variant to
use when historical tasks only partially overlap, the common case in learning-curve
and hyperparameter-transfer corpora.

\end{document}
